%\documentclass[12pt,oneside]{amsart}
\documentclass{scrartcl}

%\documentclass{article}

\usepackage{amsthm,amsmath,amssymb}
\usepackage[numbers]{natbib}

%\usepackage[colorlinks]{hyperref}
\usepackage{hyperref}
\hypersetup{
    colorlinks,%
    citecolor=black,%
    filecolor=black,%
    linkcolor=black,%
    urlcolor=black
}
\usepackage{hypernat}
\usepackage[top=2.5cm, bottom=2.5cm, left=2.8cm,
right=2.8cm]{geometry}
\usepackage{graphicx}


%\setlength{\parskip}{0pt}
%\setlength{\parsep}{0pt}
%\setlength{\headsep}{0pt}
%\setlength{\topskip}{0pt}
%\setlength{\topmargin}{0pt}
%\setlength{\topsep}{0pt}
%\setlength{\partopsep}{0pt}

\linespread{1}

%\usepackage{marvosym}

%\textwidth = 6.5 in \textheight = 9.2 in \oddsidemargin = 0.0 in
%\evensidemargin = 0.0 in \topmargin = -0.0 in \headheight = 0.0 in
%\headsep = 0.0in \parskip = 0.0in \parindent = 0.0in
%\def\baselinestretch{0.871}


\newtheorem{theorem}{Theorem}[section]
\newtheorem{proposition}[theorem]{Proposition}
\newtheorem{problem}[theorem]{Problem}
\newtheorem{fact}[theorem]{Fact}
\newtheorem{lemma}[theorem]{Lemma}
\newtheorem{defn}[theorem]{Definition}
\newtheorem{corollary}[theorem]{Corollary}
\newtheorem{remark}{Remark}[section]
\newtheorem{example}[theorem]{Example}

\newcommand{\E}{{\mathbb E}}
\newcommand{\se}{{\mathcal{E}}}
\newcommand{\R}{{\mathbb R}}
\renewcommand{\P}{{\mathbb P}}
\newcommand{\ac}{{\mathcal{A}}}
\newcommand{\A}{{\mathcal{A}}}
\newcommand{\B}{{\mathcal{B}}}
\newcommand{\C}{{\mathcal{C}}}
\newcommand{\M}{{\mathcal{M}}}
\newcommand{\Mf}{{\mathfrak{M}}}
\newcommand{\T}{{\mathcal{T}}}
\newcommand{\Z}{{\mathcal{Z}}}
\newcommand{\K}{{\mathcal{K}}}
\newcommand{\lc}{{\mathcal{L}}}
\newcommand{\Ps}{{\mathcal{P}}}
\newcommand{\G}{{\mathcal{G}}}
\newcommand{\pa}{{\mathring{p}}}
\newcommand{\F}{{\cal F}}
\newcommand{\samp}{{\mathcal{X}}}
\newcommand{\X}{{\mathcal{X}}}
\newcommand{\hi}{{\hat{\Phi}}}
\newcommand{\data}{{\text{data}}}
\newcommand{\relu}{{\text{ReLU}}}

\newcommand{\ridge}{{\hat{\beta}_{\text{ridge}}(\lambda)}}
\newcommand{\lasso}{{\hat{\beta}_{\text{lasso}}(\lambda)}}
\newcommand{\ridgemu}{{\hat{\mu}_t^{\text{ridge}}(\lambda)}}
\newcommand{\lassomu}{{\hat{\mu}_t^{\text{lasso}}(\lambda)}}


\newcounter{rcnt}[section]
\renewcommand{\thercnt}{(\roman{rcnt})}

\def\qt#1{\qquad\text{#1}}

\def\argmin{\mathop{\rm argmin}}
\def\argmax{\mathop{\rm argmax}}


\setlength{\parskip}{1.4 \medskipamount}
%\sloppy
%\linespread{1.3}

\begin{document}

\title{STAT 238 - Bayesian Statistics  \\
  Lecture Twenty Five} 
\subtitle{Spring 2026, UC Berkeley}

\author{Aditya Guntuboyina}


\date{01 April 2026}

\maketitle

\section{Markov Chains and Stationary Distributions}

Let $S$ be a state space. A sequence of random variables $\Theta^{(0)},
\Theta^{(1)}, \dots$ taking values in $S$ is a 
Markov Chain on $S$ provided the conditional distribution of  
$\Theta^{(t+1)}$ given $\Theta^{(t)} = \theta^{(t)}, \Theta^{(t-1)} =
\theta^{(t-1)}, \dots, \Theta^{(0)} = \theta^{(0)}$ only depends on
$\theta^{(t)}$. Further, if this conditional distribution is the same
for all  $t$, then the Markov chain is said to be
time-homogeneous. We will only deal with time-homogeneous Markov
Chains in this course.

A (time-homogeneous) Markov Chain is described via the conditional
distribution of $\Theta^{(t+1)}$ given $\Theta^{(t)} = x$. If $S$ is
finite or countable, then these conditional distributions can be
described by the probabilities:
\begin{align*}
  P(x, y) := \P\{\Theta^{(t+1)} = y \mid \Theta^{(t)} = x\}. 
\end{align*}
Clearly $P(x, y) \geq 0$ and $\sum_{y \in S} P(x, y) = 1$ for every $x
\in S$.

If $S$ is an open subset of $\R^d$, then one usually describes the
conditional distribution of $\Theta^{(t+1)}$ given $\Theta^{(t)} = x$
via the conditional density function $P(x, y)$:
\begin{align*}
  \P\{\Theta^{(t+1)} \in A \mid \Theta^{(t)} = x\} = \int_A P(x, y)
  dy. 
\end{align*}
Here $P(x, y) \geq 0$ and $\int_S P(x, y) dy = 1$ for all $x \in S$.

\subsection{Stationary Distribution}

A probability distribution $\pi$ on $S$ is said to be the stationary
or invariant distribution of the Markov chain $\Theta^{(0)}, 
\Theta^{(1)}, \dots$ if the following holds:
\begin{align*}
  \Theta^{(t)} \sim \pi \implies \Theta^{(t+1)} \sim \pi. 
\end{align*}
If $S$ is discrete, this can be written as:
\begin{align}\label{disc_inv}
  \pi(y) = \sum_{x \in S} \pi(x) P(x, y) ~~ \text{for every $y \in
  S$}. 
\end{align}
In the continuous case, this becomes:
\begin{align}\label{cont_inv}
  \pi(y) = \int_S \pi(x) P(x, y) dx ~~ \text{for every $y \in
  S$}. 
\end{align}

\subsection{Ergodic Theorem}
Suppose $\{\Theta^{(t)} \}$ is a time-homogeneous Markov chain that is
\textbf{irreducible}. Suppose $\pi$ is the stationary distribution of
the Markov Chain (irreducibility guarantess uniqueness of the
stationary distribution). Then
\begin{align*}
\lim_{N \rightarrow \infty}  \frac{1}{N} \sum_{t=1}^N g(\Theta^{(t)})
  = \int g(\theta) d\pi(\theta)
\end{align*}
almost surely. For a reference to this result, see
\citet[Theorem 4 and the remark after Corollary
6]{RobertsRosenthal2004}. For a comprehensive book-length reference,
see \citet{MeynTweedie2009} (Chapter 17 of this book contains this
result). 

When the state space is finite, irreducibility means that for every
pair of points $x$ and $y$ in the state space, there is a positive
probability that the chain goes from $x$ to $y$. For continuous state
spaces, irreducibility is somewhat technical to define. If you are
interested, see again the aforementioned references:
\citet{RobertsRosenthal2004, MeynTweedie2009}. 

\subsection{Detailed Balance Condition}

Note that \eqref{disc_inv} is equivalent to:
\begin{align*}
  \sum_{x \in S} \pi(y) P(y, x) = \sum_{x \in S} \pi(x) P(x, y). 
\end{align*}
This is because $\sum_{x \in S} P(y, x) = 1$. Thus a
\textbf{sufficient} condition for $\pi$ being the stationary
distribution for the Markov chain given by $P$ is:
\begin{align}\label{det_bal}
  \pi(y) P(y, x) = \pi(x) P(x, y) ~~\text{for all $x, y \in S$}. 
\end{align}
Note that one gets the same condition by also arguing from
\eqref{cont_inv} (in this case, $P(x, y)$ and $P(y, x)$ appearing in
\eqref{det_bal} should be interpreted as conditional pdfs as opposed
to pmfs). 

Here is a simple example of a Markov Chain satisfying detailed
balance.

\begin{example}[Discrete Markov Chain]
  Suppose each $\Theta^{(t)}$ takes values in a finite state space $S
  = \{1, \dots, N\}$ of cardinality $N$. Consider a graph with
  vertices $S$ and degrees $d_1, \dots, d_N$. Consider the random walk
  on the graph. The transition probabilities are then given by:
  \begin{align*}
    P(i, j) =  \frac{I\{j ~\text{is a neighbour of}~ i\}}{d_i}
  \end{align*}
  The detailed balance condition then becomes:
  \begin{align*}
    \pi(i)  \frac{I\{j ~\text{is a neighbour of}~ i\}}{d_i} = \pi(j)
    \frac{I\{i ~\text{is a neighbour of}~ j\}}{d_j}
  \end{align*}
  It is then clear that this Markov chain satisfies detailed balance
  with $\pi(i) \propto d_i$. So this $\pi(\cdot)$ is stationary. If
  the graph is connected, then this Markov chain is irreducible and
  this $\pi(\cdot)$ will be the unique stationary distribution. 
\end{example}

Our next example of a Markov Chain with detailed balance is the
Metropolis-Hastings chain.

\subsection{Metropolis-Hastings}

  Suppose $\pi$ is a probability on a state space $S$. If $S$ is
  discrete, $\pi$ should be interpreted as a pmf and if $S$ is
  continuous, $\pi$ should be interpreted as a pdf. Let $Q(x, y)$ be
  some transition probability. Again if $S$ is discrete, interpret
  $Q(x, y)$ as a pmf over $y$ for each fixed $x$. If $S$ is
  continuous, interpret $Q(x, y)$ as a pdf over $y$ for each fixed
  $x$.

  We do not assume that $\pi(\cdot)$ and $Q(\cdot, \cdot)$ are related
  in any way. For example, there is no reason for $Q$ to satisfy
  detailed balance with respect to $\pi$. In other words, for $x$ and
  $y$, it may very well happen that
  \begin{align*}
    \pi(x) Q(x, y) \neq \pi(y) Q(y, x). 
  \end{align*}
  The Metropolis-Hastings uses $Q(\cdot, \cdot)$ and $\pi(\cdot)$ to
  construct a new Markov chain $\{\Theta^{(t)}\}$ which moves as
  follows. Given $\Theta^{(t)} = x$, first use $Q(x, \cdot)$ to
  generate $y$, then take $\Theta^{(t+1)}$ to be:
  \[
\Theta^{(t+1)} =
\begin{cases}
y, & \text{with probability } \min\!\left(1, \frac{\pi(y)\,Q(y,x)}{\pi(x)\,Q(x,y)}\right), \\[1em]
x, & \text{with probability } 1 - \min\!\left(1, \frac{\pi(y)\,Q(y,x)}{\pi(x)\,Q(x,y)}\right).
\end{cases}
\]
The key is to realize that this new Markov chain satisfies detailed
balance with respect to $\pi$. To see this, we need to verify:
\begin{align}\label{mhdb}
  \pi(x) P(x, y) = \pi(y) P(y, x) ~~\text{ for all $x, y \in S$}. 
\end{align}
If $y = x$, there is nothing to prove. So let us assume that $y \neq
x$. In this case, it is easy to see that
\begin{align*}
  P(x, y) = Q(x, y) \min\!\left(1,
  \frac{\pi(y)\,Q(y,x)}{\pi(x)\,Q(x,y)}\right). 
\end{align*}
Consequently
\begin{align*}
  \pi(x) P(x, y) = \pi(x) Q(x, y) \min\!\left(1,
  \frac{\pi(y)\,Q(y,x)}{\pi(x)\,Q(x,y)}\right) =
  \min\!\left(\pi(x)\,Q(x,y), \pi(y) Q(y, x)\right). 
\end{align*}
Similarly $\pi(y) P(y, x) =   \min\!\left(\pi(x)\,Q(x,y), \pi(y) Q(y,
  x)\right)$, which proves \eqref{mhdb}.

The Metropolis-Hastings Algorithm can be interpreted as a device for
converting any Markov Chain into one which satisfies detailed balance
with respect to $\pi$. It simply adds an acceptance step to the
Markov Chain given by $Q(\cdot, \cdot)$ with acceptance probability:
\begin{align*}
  \min \left(1, \frac{\pi(y) Q(y, x)}{\pi(x) Q(x, y)} \right). 
\end{align*}
A simple special case of the Metropolis-Hastings arises when $Q(\cdot,
\cdot)$ is symmetric in the sense that $Q(x, y) = Q(y, x)$. In this
case, the acceptance probability is simply $\min(1,
\pi(y)/\pi(x))$. For example, when $Q$ is given by a random walk
Markov chain: $y = x + z$ where $z$ has a symmetric distribution
around 0 such as $z \sim N(0, \Sigma)$ or $z \sim
\text{uniform}(\prod_{j=1}^d [-a_j, a_j])$. In this special case, the
Metropolis-Hastings chain is simply referred to as the Metropolis
Chain.

\subsection{The Gibbs Sampler}

Our third example of a Markov Chain satisfying detailed balance with
respect to $\pi$ is the Gibbs sampler. Suppose the state space $S$
consists of bivariate pairs $(x_1, x_2)$. Suppose $\pi(\cdot)$ is such
that its conditionals $\pi_{1 \mid 2}(\cdot \mid x_2)$ and $\pi_{2
  \mid 1} (\cdot \mid x_1)$ are easy to sample from. Here $\pi_{1 \mid
2}(\cdot \mid x_2)$ is the conditional distribution of $x_1$ given
$x_2$ when $(x_1, x_2) \sim \pi$. Similarly $\pi_{2 \mid 1}(\cdot \mid
x_1)$ is the conditional distribution of $x_2$ given $x_1$ when
$(x_1, x_2) \sim \pi$.

The Gibbs sampler employs the following Markov Chain. Given
$\Theta^{(t)} = (x_1, x_2)$, first sample one of the indices $1$ and
$2$ with equal probability $0.5$. If we get $1$, then use the update:
\begin{align*}
  (x_1, x_2) \mapsto (x_1', x_2) ~~ \text{ where $x_1' \sim \pi_{1
  \mid 2}(\cdot
  \mid x_2)$}. 
\end{align*}
If we get $2$, then use the update:
\begin{align*}
  (x_1, x_2) \mapsto (x_1, x_2') ~~ \text{ where $x_2' \sim \pi_{2
  \mid 1}(\cdot
  \mid x_1)$}. 
\end{align*}
It turns out that this chain satisfies detailed balance with respect
to $\pi$. For this, we need to prove that
\begin{align*}
  \pi(x_1, x_2) P((x_1, x_2), (x_1', x_2')) =   \pi(x_1', x_2')
  P((x_1', x_2'), (x_1, x_2)). 
\end{align*}
The transition $(x_1, x_2) \rightarrow (x_1', x_2')$ will only be
possible if one of $x_1 = x_1'$ or $x_2 = x_2'$ hold true. Let $x_1 =
x_1'$ so we need to verify: 
\begin{align}\label{gsv}
  \pi(x_1, x_2) P((x_1, x_2), (x_1, x_2')) =   \pi(x_1, x_2')
  P((x_1, x_2'), (x_1, x_2)). 
\end{align}
From the description of the Gibbs chain, it is clear that
\begin{align*}
  P((x_1, x_2), (x_1, x_2')) = \frac{1}{2} \pi_{2 \mid 1}(x_2' \mid
  x_1) = \frac{1}{2} \frac{\pi(x_1, x_2')}{\pi_1(x_1)}
\end{align*}
where $\pi_1(\cdot)$ is the marginal  distribution of $X_1$ when
$(X_1, X_2) \sim \pi$. As a result
\begin{align*}
  \pi(x_1, x_2) P((x_1, x_2), (x_1, x_2')) = \frac{1}{2}
  \frac{\pi(x_1, x_2)\pi(x_1, x_2')}{\pi_1(x_1)}.  
\end{align*}
Clearly the right hand side above is symmetric in $x_2$ and $x_2'$
which proves \eqref{gsv}. One can similarly verify detailed balance
when $x_2 = x_2'$.






\begin{thebibliography}{99}

\bibitem[Roberts and Rosenthal(2004)]{RobertsRosenthal2004}
Roberts, G. O. and Rosenthal, J. S. (2004).
General state space Markov chains and MCMC algorithms.
\textit{Probability Surveys}, \textbf{1}, 20--71.

\bibitem[Meyn and Tweedie(2009)]{MeynTweedie2009}
Meyn, S. P. and Tweedie, R. L. (2009).
\textit{Markov Chains and Stochastic Stability}, 2nd ed.
Cambridge University Press.


  
%\bibitem[Chib and Greenberg(1995)]{chib1995}
%Chib, S. and Greenberg, E. (1995).
%Understanding the Metropolis--Hastings algorithm.
%\textit{The American Statistician}, \textbf{49}, 327--335.


%\bibitem[MacKay and Peto(1995)]{MacKay1995HierarchicalDirichlet}
%D.~J.~C. MacKay and L.~C. Peto,
%\newblock ``A hierarchical Dirichlet language model,''
%\newblock \emph{Natural Language Engineering}, vol.~1,
%no.~3, pp.~289--308, 1995.

%\bibitem[Rubin(1981)]{Rubin1981BayesianBootstrap}
%D.~B. Rubin,
%\newblock ``The Bayesian bootstrap,''
%\newblock \emph{The Annals of Statistics}, vol.~9,
%no.~1, pp.~130--134, 1981.
  
%\bibitem[Fisher(1934)]{Fisher1934}
%R.~A. Fisher,
%\newblock ``Probability, likelihood and quantity of information in the logic of
%uncertain inference,''
%\newblock \emph{Proceedings of the Royal Society of London. Series~A,
%Containing Papers of a Mathematical and Physical Character}, vol.~146,
%no.~856, pp.~1--8, 1934.

%\bibitem[Efron(2010)]{Efron}
%Efron, B. (2010)
%\textit{Large-scale inference: empirical Bayes methods for estimation,
%testing, and prediction}. Cambridge University Press, 2010.

%\bibitem[Shumway and Stoffer(2017)]{ShumwayStoffer}
%Shumway, R. H. and Stoffer, D. S. (2017)
%\textit{Time Series Analysis and Its Applications: With R
%  Examples}. Springer, 4th edition. 
  
%     \bibitem[Jaynes(2003)]{Jaynes} Jaynes, E. T, (2003)
%  \textit{Probability theory: the logic of science}. Cambridge
%  University Press, 2003.
 
%  \bibitem[Sinai(1992)]{Sinai} Sinai, Y. G, (1992)
%  \textit{Probability theory: an introductory course}. Springer
%  Verlag, 1992.  

%\bibitem[{deGroot(1986)}]{DeGrootBlackwell}DeGroot, M. H. (1986)
%       A conversation with David Blackwell.
%\newblock \emph{Statistical Science}, \textbf{1}, 40--53.

%         \bibitem[Mosteller(1987)]{Mosteller} Mosteller, F, (1987)
%  \textit{Fifty challenging problems in probability with
%    solutions}. Courier Corporation, 1987.

%    \bibitem[MacKay(2003)]{MacKay} MacKay, D. J. C., (2003)
%  \textit{Information theory, inference, and learning
%  algorithms}. Cambridge University Press, 2003.

%\bibitem[{Lindley(2013)}]{Lindley}Lindley, D. V. (2013)
%   \textit{Understanding Uncertainty}. John Wiley and sons, 2013.


\end{thebibliography}







\end{document}

%%% Local Variables:
%%% mode: latex
%%% TeX-master: t
%%% End:
